\begin{register}{H}{GPIO\_STRAP\_REG}{0x0000}\label{regdesc:GPIOSTRAPREG}
\regfield{(reserved)}{16}{16}{0000000000000000}%
\regfield{GPIO\_STRAPPING}{16}{0}{{0x00}}%
\reglabel{Reset}\regnewline%
\begin{regdesc}\begin{reglist}
\label{fielddesc:GPIOSTRAPPING}\item [GPIO\_STRAPPING] 表示 strapping 管脚的值。\\
    - Bit[0]：无效\\
    - Bit[1]: 无效\\
    - Bit[2]: GPIO8\\
    - Bit[3]: GPIO9\\
    - Bit[4]: GPIO7\\
    - Bit[5]$\sim$bit[15]：无效\\
更多有关 strapping 管脚控制的功能，见章节 \ref{mod:bootctrl} \textit{\nameref{mod:bootctrl}}。\\
 (RO)
\end{reglist}\end{regdesc}
\end{register}


\begin{register}{H}{GPIO\_OUT\_REG}{0x0004}\label{regdesc:GPIOOUTREG}
\regfield{GPIO\_OUT\_DATA\_ORIG}{32}{0}{{0x000000}}%
\reglabel{Reset}\regnewline%
\begin{regdesc}\begin{reglist}
\label{fielddesc:GPIOOUTDATAORIG}\item [GPIO\_OUT\_DATA\_ORIG] 配置简单输出模式下 GPIO0$\sim$GPIO13 和 GPIO22$\sim$GPIO29 的输出值。\\
0：低电平\\
1：高电平\\ 
bit[0]$\sim$bit[13] 和 bit[22]$\sim$bit[29] 的值分别对应 GPIO0$\sim$GPIO13 和 GPIO22$\sim$GPIO29 的输出值。\\
(R/W/SC/WTC)
\end{reglist}\end{regdesc}
\end{register}


\begin{register}{H}{GPIO\_OUT\_W1TS\_REG}{0x0008}\label{regdesc:GPIOOUTW1TSREG}
\regfield{GPIO\_OUT\_W1TS}{32}{0}{{0x000000}}%
\reglabel{Reset}\regnewline%
\begin{regdesc}\begin{reglist}
\label{fielddesc:GPIOOUTW1TS}\item [GPIO\_OUT\_W1TS] 配置是否置位 GPIO0$\sim$GPIO13 和 GPIO22$\sim$GPIO29 的输出寄存器 \hyperref[regdesc:GPIOOUTREG]{GPIO\_OUT\_REG}。\\
0：不置位\\
1：\hyperref[regdesc:GPIOOUTREG]{GPIO\_OUT\_REG} 中相应位也将置 1\\
bit[0]$\sim$bit[13] 和 bit[22]$\sim$bit[29] 分别对应 GPIO0$\sim$GPIO13 和 GPIO22$\sim$GPIO29。\\
推荐使用此寄存器来置位 \hyperref[regdesc:GPIOOUTREG]{GPIO\_OUT\_REG}。 \\ (WT)
\end{reglist}\end{regdesc}
\end{register}


\begin{register}{H}{GPIO\_OUT\_W1TC\_REG}{0x000C}\label{regdesc:GPIOOUTW1TCREG}
\regfield{GPIO\_OUT\_W1TC}{32}{0}{{0x000000}}%
\reglabel{Reset}\regnewline%
\begin{regdesc}\begin{reglist}
\label{fielddesc:GPIOOUTW1TC}\item [GPIO\_OUT\_W1TC] 配置是否清零 GPIO0$\sim$GPIO13 和 GPIO22$\sim$GPIO29 的输出寄存器 \hyperref[regdesc:GPIOOUTREG]{GPIO\_OUT\_REG}。\\ 
0：不清零\\
1：\hyperref[regdesc:GPIOOUTREG]{GPIO\_OUT\_REG} 中相应位将被清零\\
bit[0]$\sim$bit[13] 和 bit[22]$\sim$bit[29] 分别对应 GPIO0$\sim$GPIO13 和 GPIO22$\sim$GPIO29。\\
推荐使用此寄存器来清零 \hyperref[regdesc:GPIOOUTREG]{GPIO\_OUT\_REG}。 \\ (WT)
\end{reglist}\end{regdesc}
\end{register}


\begin{register}{H}{GPIO\_ENABLE\_REG}{0x0034}\label{regdesc:GPIOENABLEREG}
\regfield{GPIO\_ENABLE\_DATA}{32}{0}{{0x000000}}%
\reglabel{Reset}\regnewline%
\begin{regdesc}\begin{reglist}
\label{fielddesc:GPIOENABLEDATA}\item [GPIO\_ENABLE\_DATA] 配置是否使能 GPIO0$\sim$GPIO13 和 GPIO22$\sim$GPIO29 输出。\\
0：不使能\\
1：使能\\ 
bit[0]$\sim$bit[13] 和 bit[22]$\sim$bit[29] 分别对应 GPIO0$\sim$GPIO13 和 GPIO22$\sim$GPIO29。\\(R/W/WTC)
\end{reglist}\end{regdesc}
\end{register}


\begin{register}{H}{GPIO\_ENABLE\_W1TS\_REG}{0x0038}\label{regdesc:GPIOENABLEW1TSREG}
\regfield{GPIO\_ENABLE\_W1TS}{32}{0}{{0x000000}}%
\reglabel{Reset}\regnewline%
\begin{regdesc}\begin{reglist}
\label{fielddesc:GPIOENABLEW1TS}\item [GPIO\_ENABLE\_W1TS] 配置是否置位 GPIO0$\sim$GPIO13 和 GPIO22$\sim$GPIO29 的输出使能寄存器 \hyperref[regdesc:GPIOENABLEREG]{GPIO\_ENABLE\_REG}。\\
0：不置位\\
1：\hyperref[regdesc:GPIOENABLEREG]{GPIO\_ENABLE\_REG} 中相应位也将置 1\\
bit[0]$\sim$bit[13] 和 bit[22]$\sim$bit[29] 分别对应 GPIO0$\sim$GPIO13 和 GPIO22$\sim$GPIO29。\\
推荐使用该寄存器来置位 \hyperref[regdesc:GPIOENABLEREG]{GPIO\_ENABLE\_REG}。\\ (WT)
\end{reglist}\end{regdesc}
\end{register}


\begin{register}{H}{GPIO\_ENABLE\_W1TC\_REG}{0x003C}\label{regdesc:GPIOENABLEW1TCREG}
\regfield{GPIO\_ENABLE\_W1TC}{32}{0}{{0x000000}}%
\reglabel{Reset}\regnewline%
\begin{regdesc}\begin{reglist}
\label{fielddesc:GPIOENABLEW1TC}\item [GPIO\_ENABLE\_W1TC] 配置是否清零 GPIO0$\sim$GPIO13 和 GPIO22$\sim$GPIO29 的输出使能寄存器 \hyperref[regdesc:GPIOENABLEREG]{GPIO\_ENABLE\_REG}。\\
0：不清零\\
1：\hyperref[regdesc:GPIOENABLEREG]{GPIO\_ENABLE\_REG} 中相应位将被清零\\
bit[0]$\sim$bit[13] 和 bit[22]$\sim$bit[29] 分别对应 GPIO0$\sim$GPIO13 和 GPIO22$\sim$GPIO29。\\
推荐使用该寄存器来清零 \hyperref[regdesc:GPIOENABLEREG]{GPIO\_ENABLE\_REG}。\\ (WT)
\end{reglist}\end{regdesc}
\end{register}


\begin{register}{H}{GPIO\_IN\_REG}{0x0064}\label{regdesc:GPIOINREG}
\regfield{GPIO\_IN\_DATA\_NEXT}{32}{0}{{0x000000}}%
\reglabel{Reset}\regnewline%
\begin{regdesc}\begin{reglist}
\label{fielddesc:GPIOINDATANEXT}\item [GPIO\_IN\_DATA\_NEXT] 表示 GPIO0$\sim$GPIO13 和 GPIO22$\sim$GPIO29 的输入值。\\
每一位均代表一个管脚的输入值：\\
0：低电平\\
1：高电平\\ 
bit[0]$\sim$bit[13] 和 bit[22]$\sim$bit[29] 分别对应 GPIO0$\sim$GPIO13 和 GPIO22$\sim$GPIO29。\\
(RO)
\end{reglist}\end{regdesc}
\end{register}


\begin{register}{H}{GPIO\_STATUS\_REG}{0x0074}\label{regdesc:GPIOSTATUSREG}
\regfield{GPIO\_STATUS\_INTERRUPT}{32}{0}{{0x000000}}%
\reglabel{Reset}\regnewline%
\begin{regdesc}\begin{reglist}
\label{fielddesc:GPIOSTATUSINTERRUPT}\item [GPIO\_STATUS\_INTERRUPT] GPIO0$\sim$GPIO13 和 GPIO22$\sim$GPIO29 的中断状态寄存器，软件可配置该寄存器的值。\\
bit[0]$\sim$bit[13] 和 bit[22]$\sim$bit[29] 分别对应 GPIO0$\sim$GPIO13 和 GPIO22$\sim$GPIO29。\\
每一位均代表对应 GPIO 的中断状态：\\
    - 0：表示对应 GPIO 未产生 \hyperref[fielddesc:GPIOPINNINTTYPE]{GPIO\_PIN\regindex{n}\_INT\_TYPE} 所配置的中断，或软件配置该位寄存器的值为 0。\\
    - 1：表示对应 GPIO 产生了 \hyperref[fielddesc:GPIOPINNINTTYPE]{GPIO\_PIN\regindex{n}\_INT\_TYPE} 所配置的中断，或软件配置该位寄存器的值为 1。\\
 (R/W/WTC)
\end{reglist}\end{regdesc}
\end{register}


\begin{register}{H}{GPIO\_STATUS\_W1TS\_REG}{0x0078}\label{regdesc:GPIOSTATUSW1TSREG}
\regfield{GPIO\_STATUS\_W1TS}{32}{0}{{0x000000}}%
\reglabel{Reset}\regnewline%
\begin{regdesc}\begin{reglist}
\label{fielddesc:GPIOSTATUSW1TS}\item [GPIO\_STATUS\_W1TS] 配置是否置位 GPIO0$\sim$GPIO13 和 GPIO22$\sim$GPIO29 的中断状态寄存器 \hyperref[fielddesc:GPIOSTATUSINTERRUPT]{GPIO\_STATUS\_INTERRUPT}。\\
每一位的值可以为：\\
0：不置位\\
1：\hyperref[fielddesc:GPIOSTATUSINTERRUPT]{GPIO\_STATUS\_INTERRUPT} 中的相应位也将置 1\\
bit[0]$\sim$bit[13] 和 bit[22]$\sim$bit[29] 分别对应 GPIO0$\sim$GPIO13 和 GPIO22$\sim$GPIO29。\\
推荐使用此寄存器来置位 \hyperref[fielddesc:GPIOSTATUSINTERRUPT]{GPIO\_STATUS\_INTERRUPT}。\\
 (WT)
\end{reglist}\end{regdesc}
\end{register}


\begin{register}{H}{GPIO\_STATUS\_W1TC\_REG}{0x007C}\label{regdesc:GPIOSTATUSW1TCREG}
\regfield{GPIO\_STATUS\_W1TC}{32}{0}{{0x000000}}%
\reglabel{Reset}\regnewline%
\begin{regdesc}\begin{reglist}
\label{fielddesc:GPIOSTATUSW1TC}\item [GPIO\_STATUS\_W1TC] 配置是否清零 GPIO0$\sim$GPIO13 和 GPIO22$\sim$GPIO29 的中断状态寄存器 \hyperref[fielddesc:GPIOSTATUSINTERRUPT]{GPIO\_STATUS\_INTERRUPT}。\\
每一位的值可以为：\\
0：不清零\\
1：\hyperref[fielddesc:GPIOSTATUSINTERRUPT]{GPIO\_STATUS\_INTERRUPT} 中的相应位将被清零\\
bit[0]$\sim$bit[13] 和 bit[22]$\sim$bit[29] 分别对应 GPIO0$\sim$GPIO13 和 GPIO22$\sim$GPIO29。\\
推荐使用此寄存器来清零 \hyperref[fielddesc:GPIOSTATUSINTERRUPT]{GPIO\_STATUS\_INTERRUPT}。\\
(WT)
\end{reglist}\end{regdesc}
\end{register}


\begin{register}{H}{GPIO\_PROCPU\_INT\_REG}{0x00A4}\label{regdesc:GPIOPROCPUINTREG}
\regfield{GPIO\_PROCPU\_INT}{32}{0}{{0x000000}}%
\reglabel{Reset}\regnewline%
\begin{regdesc}\begin{reglist}
\label{fielddesc:GPIOPROCPUINT}\item [GPIO\_PROCPU\_INT] 表示 GPIO0$\sim$GPIO13 和 GPIO22$\sim$GPIO29 的 CPU 中断状态。\\ 
每位代表：\\
0：表示 CPU 中断未使能，或对应 GPIO 未产生 \hyperref[fielddesc:GPIOPINNINTTYPE]{GPIO\_PIN\regindex{n}\_INT\_TYPE} 所配置的中断。\\
1：表示 CPU 中断使能后，对应 GPIO 产生了 \hyperref[fielddesc:GPIOPINNINTTYPE]{GPIO\_PIN\regindex{n}\_INT\_TYPE} 所配置的中断。\\
bit[0]$\sim$bit[13] 和 bit[22]$\sim$bit[29] 分别对应 GPIO0$\sim$GPIO13 和 GPIO22$\sim$GPIO29。\\
如果 \hyperref[regdesc:GPIOPINNREG]{GPIO\_PIN\regindex{n}\_REG} 中 bit[13] 高电平有效，即使能 CPU 中断，则此寄存器所示的中断状态应与 \hyperref[regdesc:GPIOSTATUSREG]{GPIO\_STATUS\_REG} 中相应位的中断状态一致。 \\ (RO)
\end{reglist}\end{regdesc}
\end{register}


\begin{register}{H}{GPIO\_SDIO\_INT\_REG}{0x00A8}\label{regdesc:GPIOSDIOINTREG}
\regfield{GPIO\_SDIO\_INT}{32}{0}{{0x000000}}%
\reglabel{Reset}\regnewline%
\begin{regdesc}\begin{reglist}
\label{fielddesc:GPIOSDIOINT}\item [GPIO\_SDIO\_INT] 表示 GPIO0$\sim$GPIO13 和 GPIO22$\sim$GPIO29 的 GPIO\_SDIO\_INT 中断状态。\\ 
每位代表：\\
0：表示 GPIO\_SDIO\_INT 中断未使能，或对应 GPIO 未产生 \hyperref[fielddesc:GPIOPINNINTTYPE]{GPIO\_PIN\regindex{n}\_INT\_TYPE} 所配置的中断。\\
1：表示 GPIO\_SDIO\_INT 中断使能后，对应 GPIO 产生了 \hyperref[fielddesc:GPIOPINNINTTYPE]{GPIO\_PIN\regindex{n}\_INT\_TYPE} 所配置的中断。\\
bit[0]$\sim$bit[13] 和 bit[22]$\sim$bit[29] 分别对应 GPIO0$\sim$GPIO13 和 GPIO22$\sim$GPIO29。\\
如果 \hyperref[regdesc:GPIOPINNREG]{GPIO\_PIN\regindex{n}\_REG} 中 bit[14] 高电平有效，即使能 GPIO\_SDIO\_INT 中断，则此寄存器所示的中断状态应与 \hyperref[regdesc:GPIOSTATUSREG]{GPIO\_STATUS\_REG} 中相应位的中断状态一致。 \\ (RO)
\end{reglist}\end{regdesc}
\end{register}


\begin{register}{H}{GPIO\_STATUS\_NEXT\_REG}{0x00C4}\label{regdesc:GPIOSTATUSNEXTREG}
\regfield{GPIO\_STATUS\_INTERRUPT\_NEXT}{32}{0}{{0x000000}}%
\reglabel{Reset}\regnewline%
\begin{regdesc}\begin{reglist}
\label{fielddesc:GPIOSTATUSINTERRUPTNEXT}\item [GPIO\_STATUS\_INTERRUPT\_NEXT] 表示 GPIO0$\sim$GPIO13 和 GPIO22$\sim$GPIO29 的中断源信号。\\bit[0]$\sim$bit[13] 和 bit[22]$\sim$bit[29] 分别对应 GPIO0$\sim$GPIO13 和 GPIO22$\sim$GPIO29。\\
每位代表：\\
0：表示 GPIO 未产生 \hyperref[fielddesc:GPIOPINNINTTYPE]{GPIO\_PIN\regindex{n}\_INT\_TYPE} 所配置的中断\\
1：GPIO 产生了 \hyperref[fielddesc:GPIOPINNINTTYPE]{GPIO\_PIN\regindex{n}\_INT\_TYPE} 所配置的中断\\
上述中断可以为上升沿中断、下降沿中断、电平敏感中断或任一沿中断。\\ (RO)
\end{reglist}\end{regdesc}
\end{register}


\begin{register}{H}{GPIO\_PIN\regindex{n}\_REG (\regindex{n}: 0-13, 22-29)}{0x00D4+0x4*\regindex{n}}\label{regdesc:GPIOPINNREG}
\regfield{(reserved)}{14}{18}{00000000000000}%
\regfield{GPIO\_PIN\regindex{n}\_INT\_ENA}{5}{13}{{0x0}}%
\regfield{(reserved)}{2}{11}{00}%
\regfield{GPIO\_PIN\regindex{n}\_WAKEUP\_ENABLE}{1}{10}{{0}}%
\regfield{GPIO\_PIN\regindex{n}\_INT\_TYPE}{3}{7}{{0x0}}%
\regfield{(reserved)}{2}{5}{00}%
\regfield{GPIO\_PIN\regindex{n}\_SYNC1\_BYPASS}{2}{3}{{0x0}}%
\regfield{GPIO\_PIN\regindex{n}\_PAD\_DRIVER}{1}{2}{{0}}%
\regfield{GPIO\_PIN\regindex{n}\_SYNC2\_BYPASS}{2}{0}{{0x0}}%
\reglabel{Reset}\regnewline%
\begin{regdesc}\begin{reglist}
\label{fielddesc:GPIOPINNSYNC2BYPASS}\item [GPIO\_PIN\regindex{n}\_SYNC2\_BYPASS] 配置是否使能 GPIO 输入数据在第二级同步中 HP IO MUX 运行时钟上升沿同步或下降沿同步。\\
0：不同步\\
1：在下降沿进行同步\\
2：在上升沿进行同步\\
3：在上升沿进行同步\\ (R/W)
\label{fielddesc:GPIOPINNPADDRIVER}\item [GPIO\_PIN\regindex{n}\_PAD\_DRIVER] 配置选择管脚的输出驱动模式。\\
0：正常输出\\
1：开漏输出 \\ (R/W)
\label{fielddesc:GPIOPINNSYNC1BYPASS}\item [GPIO\_PIN\regindex{n}\_SYNC1\_BYPASS] 配置是否使能 GPIO 输入数据在第一级同步中 HP IO MUX 运行时钟上升沿同步或下降沿同步。\\
0：不同步\\
1：在下降沿进行同步\\
2：在上升沿进行同步\\
3：在上升沿进行同步\\ (R/W)
\label{fielddesc:GPIOPINNINTTYPE}\item [GPIO\_PIN\regindex{n}\_INT\_TYPE] 配置 GPIO 中断类型。\\
0：关闭 GPIO 中断\\
1：上升沿触发\\
2：下降沿触发\\
3：任一沿触发\\
4：低电平触发\\
5：高电平触发\\ (R/W)
\label{fielddesc:GPIOPINNWAKEUPENABLE}\item [GPIO\_PIN\regindex{n}\_WAKEUP\_ENABLE] 配置是否使能 GPIO 唤醒功能。\\
0：不使能\\
1：使能\\
该功能只能将 CPU 从 Light-sleep 中唤醒。 \\ (R/W)
\item [见下页]
\end{reglist}\end{regdesc}
\end{register}

\addtocounter{Regfloat}{-1}
\begin{register}{H}{GPIO\_PIN\regindex{n}\_REG (\regindex{n}: 0-29)}{0x00D4+0x4*\regindex{n}}
\begin{regdesc}\begin{reglist}
\item [接上页]
\label{fielddesc:GPIOPINNINTENA}\item [GPIO\_PIN\regindex{n}\_INT\_ENA] 配置是否使能 GPIO\regindex{n} 产生 CPU 中断\\
- bit[13]：配置是否使能 GPIO\_EXT\_REG\_INT 中断：\\
    0：不使能\\
    1：使能\\
- bit[15]：配置是否使能 GPIO\_SDIO\_INT 中断：\\
    0：不使能\\
    1：使能\\
- bit[14]、bit[16]$\sim$bit[17]：无效\\
 (R/W)
\end{reglist}\end{regdesc}
\end{register}

\begin{register}{H}{GPIO\_FUNC\regindex{m}\_IN\_SEL\_CFG\_REG (\regindex{m}: 6-17)}{0x02D4+0x4*(\regindex{m}-6)}\label{regdesc:GPIOFUNCMINSELCFGREG}
\regfield{(reserved)}{23}{9}{00000000000000000000000}%
\regfield{GPIO\_SIG\regindex{m}\_IN\_SEL}{1}{8}{{0}}%
\regfield{GPIO\_FUNC\regindex{m}\_IN\_INV\_SEL}{1}{7}{{0}}%
\regfield{GPIO\_FUNC\regindex{m}\_IN\_SEL}{7}{0}{{0x60}}%
\reglabel{Reset}\regnewline%
\begin{regdesc}\begin{reglist}
\label{fielddesc:GPIOFUNCMINSEL}\item [GPIO\_FUNC\regindex{m}\_IN\_SEL] 配置从 22 个 GPIO 中选择一个管脚连接输入信号 \regindex{m}。\\
0x0：选择 GPIO0\\
0x1：选择 GPIO1\\
......\\
0x1B：选择 GPIO27\\
0x1C：选择 GPIO28\\
或者\\
0x40：选择恒高电平输入\\
0x60：选择恒低电平输入\\ (R/W)
\label{fielddesc:GPIOFUNCNININVSEL}\item [GPIO\_FUNC\regindex{m}\_IN\_INV\_SEL] 配置是否反转输入值。\\
0：不反转\\
1：反转\\ (R/W)
\label{fielddesc:GPIOSIGNINSEL}\item [GPIO\_SIG\regindex{m}\_IN\_SEL] 配置是否通过 HP GPIO 交换矩阵连接信号。\\
0：旁路 HP GPIO 交换矩阵，即直接通过 HP IO MUX 连接信号与外设\\
1：通过 HP GPIO 交换矩阵连接信号与外设\\
 (R/W)
\end{reglist}\end{regdesc}
\end{register}

\begin{register}{H}{GPIO\_FUNC\regindex{p}\_IN\_SEL\_CFG\_REG (\regindex{p}: 27-35)}{0x0340+0x4*(\regindex{p}-27)}\label{regdesc:GPIOFUNCPINSELCFGREG}
\regfield{(reserved)}{23}{9}{00000000000000000000000}%
\regfield{GPIO\_SIG\regindex{p}\_IN\_SEL}{1}{8}{{0}}%
\regfield{GPIO\_FUNC\regindex{p}\_IN\_INV\_SEL}{1}{7}{{0}}%
\regfield{GPIO\_FUNC\regindex{p}\_IN\_SEL}{7}{0}{{0x60}}%
\reglabel{Reset}\regnewline%
\begin{regdesc}\begin{reglist}
\label{fielddesc:GPIOFUNCPINSEL}\item [GPIO\_FUNC\regindex{p}\_IN\_SEL] 配置从 22 个 GPIO 中选择一个管脚连接输入信号 \regindex{p}。\\
0x0：选择 GPIO0\\
0x1：选择 GPIO1\\
......\\
0x1B：选择 GPIO27\\
0x1C：选择 GPIO28\\
Or\\
0x40：选择恒高电平输入\\
0x60：选择恒低电平输入\\ (R/W)
\label{fielddesc:GPIOFUNCPININVSEL}\item [GPIO\_FUNC\regindex{p}\_IN\_INV\_SEL]  配置是否反转输入值。\\
0：不反转\\
1：反转\\ (R/W)
\label{fielddesc:GPIOSIGPINSEL}\item [GPIO\_SIG\regindex{p}\_IN\_SEL] 配置是否通过 HP GPIO 交换矩阵连接信号。\\
0：旁路 HP GPIO 交换矩阵，即直接通过 HP IO MUX 连接信号与外设\\
1：通过 HP GPIO 交换矩阵连接信号与外设\\
 (R/W)
\end{reglist}\end{regdesc}
\end{register}

\begin{register}{H}{GPIO\_FUNC\regindex{q}\_IN\_SEL\_CFG\_REG (\regindex{q}: 46-47)}{0x038C+0x4*(\regindex{q}-46)}\label{regdesc:GPIOFUNCQINSELCFGREG}
\regfield{(reserved)}{23}{9}{00000000000000000000000}%
\regfield{GPIO\_SIG\regindex{q}\_IN\_SEL}{1}{8}{{0}}%
\regfield{GPIO\_FUNC\regindex{q}\_IN\_INV\_SEL}{1}{7}{{0}}%
\regfield{GPIO\_FUNC\regindex{q}\_IN\_SEL}{7}{0}{{0x60}}%
\reglabel{Reset}\regnewline%
\begin{regdesc}\begin{reglist}
\label{fielddesc:GPIOFUNCQINSEL}\item [GPIO\_FUNC\regindex{q}\_IN\_SEL] 配置从 22 个 GPIO 中选择一个管脚连接输入信号 \regindex{q}。\\
0x0：选择 GPIO0\\
0x1：选择 GPIO1\\
......\\
0x1B：选择 GPIO27\\
0x1C：选择 GPIO28\\
Or\\
0x40：选择恒高电平输入\\
0x60：选择恒低电平输入\\ (R/W)
\label{fielddesc:GPIOFUNCQININVSEL}\item [GPIO\_FUNC\regindex{q}\_IN\_INV\_SEL]  配置是否反转输入值。\\
0：不反转\\
1：反转\\ (R/W)
\label{fielddesc:GPIOSIGQINSEL}\item [GPIO\_SIG\regindex{q}\_IN\_SEL] 配置是否通过 HP GPIO 交换矩阵连接信号。\\
0：旁路 HP GPIO 交换矩阵，即直接通过 HP IO MUX 连接信号与外设\\
1：通过 HP GPIO 交换矩阵连接信号与外设\\
 (R/W)
\end{reglist}\end{regdesc}
\end{register}

\begin{register}{H}{GPIO\_FUNC\regindex{r}\_IN\_SEL\_CFG\_REG (\regindex{r}: 64-69)}{0x03D4+0x4*(\regindex{r}-64)}\label{regdesc:GPIOFUNCRINSELCFGREG}
\regfield{(reserved)}{23}{9}{00000000000000000000000}%
\regfield{GPIO\_SIG\regindex{r}\_IN\_SEL}{1}{8}{{0}}%
\regfield{GPIO\_FUNC\regindex{r}\_IN\_INV\_SEL}{1}{7}{{0}}%
\regfield{GPIO\_FUNC\regindex{r}\_IN\_SEL}{7}{0}{{0x60}}%
\reglabel{Reset}\regnewline%
\begin{regdesc}\begin{reglist}
\label{fielddesc:GPIOFUNCRINSEL}\item [GPIO\_FUNC\regindex{r}\_IN\_SEL] 配置从 22 个 GPIO 中选择一个管脚连接输入信号 \regindex{r}。\\
0x0：选择 GPIO0\\
0x1：选择 GPIO1\\
......\\
0x1B：选择 GPIO27\\
0x1C：选择 GPIO28\\
Or\\
0x40：选择恒高电平输入\\
0x60：选择恒低电平输入\\ (R/W)
\label{fielddesc:GPIOFUNCRININVSEL}\item [GPIO\_FUNC\regindex{r}\_IN\_INV\_SEL]  配置是否反转输入值。\\
0：不反转\\
1：反转\\ (R/W)
\label{fielddesc:GPIOSIGRINSEL}\item [GPIO\_SIG\regindex{r}\_IN\_SEL] 配置是否通过 HP GPIO 交换矩阵连接信号。\\
0：旁路 HP GPIO 交换矩阵，即直接通过 HP IO MUX 连接信号与外设\\
1：通过 HP GPIO 交换矩阵连接信号与外设\\
 (R/W)
\end{reglist}\end{regdesc}
\end{register}

\begin{register}{H}{GPIO\_FUNC\regindex{t}\_IN\_SEL\_CFG\_REG (\regindex{t}: 72-74)}{0x03F4+0x4*(\regindex{t}-72)}\label{regdesc:GPIOFUNCTINSELCFGREG}
\regfield{(reserved)}{23}{9}{00000000000000000000000}%
\regfield{GPIO\_SIG\regindex{t}\_IN\_SEL}{1}{8}{{0}}%
\regfield{GPIO\_FUNC\regindex{t}\_IN\_INV\_SEL}{1}{7}{{0}}%
\regfield{GPIO\_FUNC\regindex{t}\_IN\_SEL}{7}{0}{{0x60}}%
\reglabel{Reset}\regnewline%
\begin{regdesc}\begin{reglist}
\label{fielddesc:GPIOFUNCTINSEL}\item [GPIO\_FUNC\regindex{t}\_IN\_SEL] 配置从 22 个 GPIO 中选择一个管脚连接输入信号 \regindex{t}。\\
0x0：选择 GPIO0\\
0x1：选择 GPIO1\\
......\\
0x1B：选择 GPIO27\\
0x1C：选择 GPIO28\\
Or\\
0x40：选择恒高电平输入\\
0x60：选择恒低电平输入\\ (R/W)
\label{fielddesc:GPIOFUNCTININVSEL}\item [GPIO\_FUNC\regindex{t}\_IN\_INV\_SEL]  配置是否反转输入值。\\
0：不反转\\
1：反转\\ (R/W)
\label{fielddesc:GPIOSIGTINSEL}\item [GPIO\_SIG\regindex{t}\_IN\_SEL] 配置是否通过 HP GPIO 交换矩阵连接信号。\\
0：旁路 HP GPIO 交换矩阵，即直接通过 HP IO MUX 连接信号与外设\\
1：通过 HP GPIO 交换矩阵连接信号与外设\\
 (R/W)
\end{reglist}\end{regdesc}
\end{register}

\begin{register}{H}{GPIO\_FUNC\regindex{u}\_IN\_SEL\_CFG\_REG (\regindex{u}: 97-100)}{0x0458+0x4*(\regindex{u}-97)}\label{regdesc:GPIOFUNCUINSELCFGREG}
\regfield{(reserved)}{23}{9}{00000000000000000000000}%
\regfield{GPIO\_SIG\regindex{u}\_IN\_SEL}{1}{8}{{0}}%
\regfield{GPIO\_FUNC\regindex{u}\_IN\_INV\_SEL}{1}{7}{{0}}%
\regfield{GPIO\_FUNC\regindex{u}\_IN\_SEL}{7}{0}{{0x60}}%
\reglabel{Reset}\regnewline%
\begin{regdesc}\begin{reglist}
\label{fielddesc:GPIOFUNCUINSEL}\item [GPIO\_FUNC\regindex{u}\_IN\_SEL] 配置从 22 个 GPIO 中选择一个管脚连接输入信号 \regindex{u}。\\
0x0：选择 GPIO0\\
0x1：选择 GPIO1\\
......\\
0x1B：选择 GPIO27\\
0x1C：选择 GPIO28\\
Or\\
0x40：选择恒高电平输入\\
0x60：选择恒低电平输入\\ (R/W)
\label{fielddesc:GPIOFUNCUININVSEL}\item [GPIO\_FUNC\regindex{u}\_IN\_INV\_SEL]  配置是否反转输入值。\\
0：不反转\\
1：反转\\ (R/W)
\label{fielddesc:GPIOSIGUINSEL}\item [GPIO\_SIG\regindex{u}\_IN\_SEL] 配置是否通过 HP GPIO 交换矩阵连接信号。\\
0：旁路 HP GPIO 交换矩阵，即直接通过 HP IO MUX 连接信号与外设\\
1：通过 HP GPIO 交换矩阵连接信号与外设\\
 (R/W)
\end{reglist}\end{regdesc}
\end{register}

\begin{register}{H}{GPIO\_FUNC\regindex{n}\_OUT\_SEL\_CFG\_REG (\regindex{n}: 0-13, 22-29)}{0x0AD4+0x4*\regindex{n}}\label{regdesc:GPIOFUNCNOUTSELCFGREG}
\regfield{(reserved)}{20}{12}{00000000000000000000}%
\regfield{GPIO\_FUNC\regindex{n}\_OE\_INV\_SEL}{1}{11}{{0}}%
\regfield{GPIO\_FUNC\regindex{n}\_OE\_SEL}{1}{10}{{0}}%
\regfield{GPIO\_FUNC\regindex{n}\_OUT\_INV\_SEL}{1}{9}{{0}}%
\regfield{GPIO\_FUNC\regindex{n}\_OUT\_SEL}{9}{0}{{0x100}}%
\reglabel{Reset}\regnewline%
\begin{regdesc}\begin{reglist}
\label{fielddesc:GPIOFUNCNOUTSEL}\item [GPIO\_FUNC\regindex{n}\_OUT\_SEL] 配置从外设信号中选择一个信号 $Y$ (0 <= $Y$ < 256) 输出至 GPIO\regindex{n}。\\
0：选择信号 0\\
1：选择信号 1\\
......\\
254：选择信号 254\\
255：选择信号 255\\
256：\hyperref[regdesc:GPIOOUTREG]{GPIO\_OUT\_REG} 和 \hyperref[regdesc:GPIOENABLEREG]{GPIO\_ENABLE\_REG} 中的 bit[\regindex{n}] 将用作输出值和输出使能信号。\\
信号列表见表 \ref{tab:iomuxgpio-periph-signals-via-hp-gpio-matrix-cn}。\\
(R/W/SC)
\label{fielddesc:GPIOFUNCNOUTINVSEL}\item [GPIO\_FUNC\regindex{n}\_OUT\_INV\_SEL] 配置是否反转输出值。\\
0：不反转\\
1：反转\\ (R/W/SC)
\label{fielddesc:GPIOFUNCNOESEL}\item [GPIO\_FUNC\regindex{n}\_OE\_SEL] 配置选择输出使能信号。\\
0：使用外设的输出使能信号 \\ 
1：强制使用 \hyperref[regdesc:GPIOENABLEREG]{GPIO\_ENABLE\_REG} 中的 bit[\regindex{n}] 用作输出使能信号。\\ (R/W)
\label{fielddesc:GPIOFUNCNOEINVSEL}\item [GPIO\_FUNC\regindex{n}\_OE\_INV\_SEL] 配置是否反转输出使能信号。\\
0：不反转\\
1：反转\\ (R/W)
\end{reglist}\end{regdesc}
\end{register}


\begin{register}{H}{GPIO\_CLOCK\_GATE\_REG}{0x0DF8}\label{regdesc:GPIOCLOCKGATEREG}
\regfield{(reserved)}{31}{1}{0000000000000000000000000000000}%
\regfield{GPIO\_CLK\_EN}{1}{0}{{1}}%
\reglabel{Reset}\regnewline%
\begin{regdesc}\begin{reglist}
\label{fielddesc:GPIOCLKEN}\item [GPIO\_CLK\_EN] 配置是否使能时钟门控。\\
0：不使能\\
1：此位置 1，则时钟自由运转。 \\ (R/W)
\end{reglist}\end{regdesc}
\end{register}


\begin{register}{H}{GPIO\_DATE\_REG}{0x0DFC}\label{regdesc:GPIODATEREG}
\regfield{(reserved)}{4}{28}{0000}%
\regfield{GPIO\_DATE}{28}{0}{{0x2412290}}%
\reglabel{Reset}\regnewline%
\begin{regdesc}\begin{reglist}
\label{fielddesc:GPIODATE}\item [GPIO\_DATE] 版本控制寄存器。 \\ (R/W)
\end{reglist}\end{regdesc}
\end{register}
